\subsection*{53. Sales Talent Is About Buyer Diagnosis and Agreement Quality}

A salesperson's output is not merely the number of people who say yes. It is the production of \textbf{qualified commitments that the business can fulfill profitably and that the buyer understands}. Great performers discover the customer's problem, identify the decision structure, test value, handle objections without pressure tactics, advance the process, negotiate sustainable terms, and hand over accurate expectations. Revenue without appropriate fit may later become refunds, churn, reputation loss, or implementation debt.

Chapter 21 described the sales process; the hiring question is whether a person can execute it in the target buying environment. D2C admissions guidance, college partnerships, enterprise consultative selling, and transactional retail selling demand different expertise. Avoid selecting an enterprise salesperson solely because they once owned a large quota at a mature brand; the relevant question is how much of their success they created versus inherited.

\subsection*{54. Sales Role-Play: Let the Customer Be a Poor Fit}

A high-value sales simulation should include an \textbf{unsellable or inappropriate buyer}. For an educational product, the role-player says they need a guaranteed placement within 3 months and cannot commit weekly study hours. The candidate's incentives are ostensibly to close enrollment. Observe whether they clarify goals and constraints, explain what the offering can and cannot do, offer a suitable alternative if 1 exists, and decline the sale when the promise cannot be responsibly made.

This differentiates professional sales judgment from short-term persuasion. A candidate who extracts a deposit through urgency tactics may score well on a naive close-rate metric while creating institutional harm. In enterprise sales, the analogous test is an unauthorized budget owner or a requested service-level guarantee beyond delivery capacity. \textbf{Responsible disqualification} is a sign of maturity.

\subsection*{55. Separate Sales Discovery from Sales Presentation}

An effective practical battery can include 3 short stages: \textbf{discovery} (understand an ambiguous buyer), \textbf{value articulation} (present an appropriately bounded proposal), and \textbf{next-step management} (secure a specific authorized action or decide not to pursue). Give the candidate consistent buyer information and evaluate which facts they elicit themselves. An excellent pitch to the wrong customer is not good selling.

Add a written exercise: summarize the account in a CRM note, identify missing stakeholders, define the next milestone, and estimate why it might fail. Strong notes distinguish confirmed information from guesses. If the fictional buying company requires legal, IT, and finance review, a salesperson who predicts closing tomorrow because the user likes the demo has not mapped the deal.

\subsection*{56. Sales Process Memory Is Not Sales Judgment}

Candidates may recite MEDDICC, SPIN, Challenger, BANT, or a custom framework fluently. These can organize behavior but are not proof of performance. Ask the salesperson to apply the relevant concepts to an unfamiliar case, prioritize which information matters, and explain their strategy when stakeholders disagree. A person who memorizes a qualification acronym but cannot recognize that there is no funded business problem is not qualified by the jargon.

Probe a real opportunity from previous employment: what was the initial trigger, which stakeholders had power, what did the candidate personally do, what was the alternative, what concessions were made, who delivered the outcome, and what happened later? Compare the account with a reference where authorized. The objective is contribution and causal plausibility, not cross-examining the candidate into a predetermined story.

\subsection*{57. Distinguish Inherited Pipeline from Created Pipeline}

A salesperson who closed \text{₹}2~crore of deals may have inherited warm inbound leads, established brand credibility, discount authority, and a mature technical team. Another who closed \text{₹}40~lakh through founder-like outbound selling to skeptical buyers may have demonstrated stronger fit for a small venture. Neither number can be interpreted without \textbf{sales cycle, account source, average contract value, quota design, gross margin, discounting, churn, team support, and territory quality}.

Ask for an anonymized deal decomposition and focus on the candidate's controllable choices. When the information cannot be verified, treat the claim as weak evidence. For early founder-led companies, experience selling without a large brand and willingness to perform honest discovery may matter more than senior title. A mature enterprise account executive could still be excellent; test the context rather than stereotyping experience.

\subsection*{58. The Salesperson Scorecard}

For a founder-assisted B2B seller, a provisional rubric might be:
\begin{itemize}[leftmargin=*]
\item Diagnosis and qualification: 25\%
\item Stakeholder mapping: 20\%
\item Commercial reasoning: 15\%
\item Credible value communication: 15\%
\item Process and follow-through: 15\%
\item Ethical judgment: 10\%
\end{itemize}
Set non-negotiable gates for misrepresentation, unauthorized commitments, and confidentiality. Add a practical task and 2 references before considering production targets.

For a D2C admissions adviser, increase weighting for listening, clarity, suitability assessment, and respectful handling of financial constraints. For enterprise sales, increase discovery, account planning, procurement, contracting, and complex negotiation. For partner sales, test ability to understand partner economics and rules of engagement. 1 scorecard cannot be carried unchanged from consumer education to global enterprise software.

\subsection*{59. Evaluate Negotiation Through a Structured Trade-Off}

A negotiation case should not reward whoever speaks louder or holds the price longest. Give the candidate a customer asking for a discount, additional support, an accelerated launch, and unfavorable payment terms. Provide confidential seller-side capacity and margin constraints. Ask the candidate to identify what value each concession creates, which concessions are costly, which are reversible, and what can be exchanged for an enforceable commitment.

Good judgment includes protecting delivery capacity, documenting promises, and walking away when the economics fail. Have the candidate draft a short follow-up email confirming scope, timeline, approval rights, and next steps. Often the written follow-up reveals more about commercial discipline than the performance of the negotiation itself.

\subsection*{60. Hire Customer-Success Talent for Value Realization}

Customer-success specialists occupy a boundary between operations, education, support, retention, and commercial expansion. The best are not necessarily the most agreeable people. They must recognize early failure signals, diagnose poor adoption, coordinate recovery, set boundaries, and help customers realize the value already purchased. Their work is judged by outcomes and sustainable relationships, not simply friendliness or volume of check-in meetings.

Use a scenario where a paying customer is disengaged and blames the product. Ask the candidate what data they need, how they would speak to the customer, where they would take responsibility, and how they would avoid promising an impossible fix. In education, assess capability to support students without doing their assignments for them. In enterprise settings, test stakeholder communication, implementation milestones, and the ability to identify a poor-fit account rather than suppress churn through manipulation.

\subsection*{61. A Side-by-Side Comparison: Consultant, Operator, Salesperson}

A comparative summary of operational and commercial roles is given in Table~\ref{tab:commercial_role_comparison}.

\begin{table*}[t]
\centering
\small
\caption{Comparative Selection Profiles for Commercial and Operational Roles}
\label{tab:commercial_role_comparison}
\begin{tabularx}{\textwidth}{p{2.5cm} p{4.2cm} p{4.5cm} X}
\toprule
\textbf{Job} & \textbf{Core Decision} & \textbf{Best Short Sample} & \textbf{Most Dangerous False Positive} \\
\midrule
Consultant & What problem is worth solving and how? & Ambiguous client case with defensible recommendation & Polished story without diagnosis \\
Operator & What must happen next when plans collide? & Timed in-tray and incident response & Heroic improviser without process control \\
Salesperson & Is there an authorized, beneficial exchange? & Live buyer discovery and qualification & Charismatic closer who mis-sells \\
Customer-success lead & How can promised value be realized and recovered? & At-risk customer diagnosis & Pleasant account manager who hides problems \\
\bottomrule
\end{tabularx}
\end{table*}

These roles overlap, especially in tiny firms, but selection should reflect the expected distribution of work. A founder may need to perform all 4 until there is enough repeatable demand to justify specialization. The right moment to hire is when a bounded job and economic case exist, not when the founder is tired of switching among roles.
